\documentclass[11pt]{article}
\usepackage[a4paper,margin=25mm]{geometry}
\usepackage{amsmath,amssymb,mathtools}
\usepackage{fontspec,unicode-math}
\setmainfont{Latin Modern Roman}
\setsansfont{Latin Modern Sans}
\setmathfont{Latin Modern Math}
\usepackage{xurl,hyperref}
\hypersetup{colorlinks=true,linkcolor=blue,urlcolor=blue,pdftitle={Problem statement}}
\newcommand{\sref}[2]{\hyperlink{source:#1:#2}{[S#2]}}
\newcommand{\cataloguescope}[1]{\paragraph{Mathematical question.}#1}
\begin{document}
\section{Dynamics of exponential-trigonometric entire maps}\label{statement}

\subsection{Definitions and mathematical statement}

\paragraph{Shared notation.}$\mathbb N=\{1,2,\ldots\}$, and $\mathbb Z,\mathbb Q,\mathbb R,\mathbb C$ denote the integers, rationals, reals and complex numbers. Any other conventions are specified locally.


A transcendental entire map is a holomorphic function $f:\mathbb C\to\mathbb C$ that is not a polynomial. Its dynamics means the iterates $f^n$, their orbit behaviour, and the stable (Fatou) and Julia sets: the Fatou set is the largest open set where the iterates form a normal family, and its complement is $J(f)$.
\paragraph{Statement recorded in the supplied source.}
For complex parameters $\lambda$, describe the dynamics of both $f_\lambda(z)=\lambda e^z\sin z$ and $g_\lambda(z)=\lambda e^z\cos z$. \sref{statement}{2}
\subsection{Short English statement}
Describe the iteration dynamics of the exponential-times-sine and exponential-times-cosine families as their complex parameter varies.
\subsection{Sources}
\begin{description}
\item[\hypertarget{source:statement:1}{S1}] ULAM.AI and the UnsolvedMath Contributors, UnsolvedMath: A Curated Collection of Open Mathematics Problems, catalogue entry. CC BY 4.0. \url{https://huggingface.co/datasets/ulamai/UnsolvedMath}.
\item[\hypertarget{source:statement:2}{S2}] Robert Devaney, ``Open Questions in Non-Rational Complex Dynamics'', in Ben Bielefeld and Mikhail Lyubich (eds.), Problems in Holomorphic Dynamics (1992), second problem in Section 1, pp. 35; arXiv:math/9205209v1. \url{https://arxiv.org/pdf/math/9205209}.
\end{description}
\label{end:statement}
\end{document}
